\begin{ouverture}[translation]{13}{Translations et parallélogrammes}{Partie 4 · Géométrie et espace}
{Glissement · Conservations · Parallélogrammes · Frises et pavages}
\objectif{Construire l’image d’une figure par une translation.}
\objectif{Utiliser les propriétés de conservation et le lien avec le parallélogramme.}
\objectif{Reconnaître symétries, demi-tours, translations et parallélogrammes particuliers.}
\prerequis{Symétrie axiale, demi-tour et parallélogrammes (Cinquième).}
\end{ouverture}
\pagestyle{otMain}\essentiel\fontsize{11}{15}\selectfont
\phantomsection\label{cours-13-1}
\essSection{1}{Faire glisser une figure}
\begin{definition}La translation qui transforme A en A’ fait glisser toute figure selon la direction de la droite $(AA')$, dans le sens de A vers A’, d’une longueur égale à $AA'$. Elle déplace tous les points de la même façon, sans faire tourner la figure.\end{definition}
\label{fig-cours-13-1-cours-translation}\figureGeo{cours-translation}

\begin{propriete}Une translation conserve les longueurs, les angles, l’alignement, le parallélisme et les aires. L’image d’une droite est une droite parallèle. Si A a pour image A’ et B pour image B’, alors $[AB]$ et $[A'B']$ sont parallèles et de même longueur, et ABB’A’ est un parallélogramme, éventuellement aplati.\end{propriete}
\begin{savoirfaire}[Construire une image]Sur quadrillage, compter le déplacement de A vers A’ et le reproduire depuis chaque sommet. Sans quadrillage, construire chaque image B’ comme quatrième sommet du parallélogramme ABB’A’, à la règle et au compas, puis relier les images dans le même ordre.\end{savoirfaire}
\clearpage
\phantomsection\label{cours-13-2}
\essSection{2}{Reconnaître une transformation}
\begin{definition}[Trois transformations]Une symétrie axiale retourne la figure par rapport à une droite. Un demi-tour, ou symétrie centrale, la fait tourner d’un demi-tour autour d’un point. Une translation la fait glisser sans la tourner. Toutes trois conservent les longueurs, les angles et les aires.\end{definition}
\begin{savoirfaire}[Distinguer]Par une translation, un segment et son image sont parallèles, de même longueur et orientés dans le même sens. Par un demi-tour, ils sont parallèles mais de sens opposés. Une symétrie axiale inverse le sens de rotation dans lequel on parcourt les sommets d’un polygone.\end{savoirfaire}
\begin{propriete}[Rappels sur les parallélogrammes]Un quadrilatère non croisé dont les diagonales ont le même milieu est un parallélogramme. Un parallélogramme dont les diagonales ont la même longueur est un rectangle ; s’il a des diagonales perpendiculaires, c’est un losange ; les deux à la fois, c’est un carré.\end{propriete}
\begin{exemple}[Frises et pavages]Une frise répète un motif par une même translation ; un pavage le répète dans deux directions. Les décors de l’Alhambra de Grenade et les dessins de M. C. Escher utilisent ainsi des translations.\end{exemple}

\clearpage\rubrique{Découverte}
\begin{blocLivre}{ex-13-1}\begin{exercice}[1]{Décrire le glissement}
Dans le repère, la translation transforme P(0 ; 0) en Q(3 ; 1). Quelle est l’image du point A(1 ; 2) ?
\label{fig-ex-13-1}\figureGeo{trans-01}
\end{exercice}\end{blocLivre}
\begin{blocLivre}{ex-13-2}\begin{exercice}[1]{Construire un triangle image}
Construire l’image du triangle de sommets A(1 ; 1), B(3 ; 1) et C(2 ; 3) par le glissement de 3 carreaux vers la droite et 2 vers le haut.
\label{fig-ex-13-2}\figureGeo{trans-02}
\end{exercice}\end{blocLivre}
\begin{blocLivre}{ex-13-3}\begin{exercice}[1]{Un déplacement vers la gauche}
Une translation déplace de 4 carreaux vers la gauche et 1 vers le bas. Donner l’image du point B(2 ; 3).
\label{fig-ex-13-3}\figureGeo{trans-03}
\end{exercice}\end{blocLivre}
\begin{blocLivre}{ex-13-4}\begin{exercice}[1]{Conserver des mesures}
Le triangle ABC a $AB=5$ cm, $\widehat{BAC}=42^\circ$ et une aire de 9 cm$^2$. Donner ces mesures pour son image A’B’C’ par une translation.
\label{fig-ex-13-4}\figureGeo{trans-04}
\end{exercice}\end{blocLivre}
\begin{blocLivre}{ex-13-5}\begin{exercice}[1]{Reconnaître une transformation}
Les drapeaux 1, 2 et 3 sont des images du drapeau F. Pour chacun, dire s’il est obtenu par une symétrie axiale, un demi-tour ou une translation.
\label{fig-ex-13-5}\figureGeo{trans-05}
\end{exercice}\end{blocLivre}

\rubrique{Entraînement}
\begin{blocLivre}{ex-13-6}\begin{exercice}[2]{Un parallélogramme}
A a pour image A’ et B pour image B’ par une même translation. Les quatre points sont distincts et non alignés. Justifier que ABB’A’ est un parallélogramme.
\label{fig-ex-13-6}\figureGeo{trans-06}
\end{exercice}\end{blocLivre}
\begin{blocLivre}{ex-13-7}\begin{exercice}[2]{Reconnaître par les diagonales}
Dans chaque figure, les diagonales $[AC]$ et $[BD]$ du quadrilatère ABCD se coupent en O ; les codages sont des données. Donner la nature de chaque quadrilatère.
\label{fig-ex-13-7}\figureGeo{trans-07}
\end{exercice}\end{blocLivre}
\begin{blocLivre}{ex-13-8}\begin{exercice}[2]{Une frise}
Les débuts de cinq motifs identiques d’une frise sont espacés de 4 cm. Quel glissement mène du premier motif au suivant ? du premier au cinquième ?
\label{fig-ex-13-8}\figureGeo{trans-08}
\end{exercice}\end{blocLivre}
\begin{blocLivre}{ex-13-9}\begin{exercice}[2]{Image d’une droite}
Une translation transforme la droite $(d)$ en $(d')$ et la droite $(\Delta)$ en $(\Delta')$. Les droites $(d)$ et $(\Delta)$ forment un angle de $35^\circ$. Que peut-on dire de $(d)$ et $(d')$ ? Quel angle forment $(d')$ et $(\Delta')$ ?
\label{fig-ex-13-9}\figureGeo{trans-09}
\end{exercice}\end{blocLivre}
\begin{blocLivre}{ex-13-10}\begin{exercice}[2]{Deux déplacements}
On déplace un point de 5 carreaux vers la droite et 2 vers le haut, puis de 3 vers la gauche et 4 vers le haut. Quel glissement unique donne le même résultat ?
\label{fig-ex-13-10}\figureGeo{trans-10}
\end{exercice}\end{blocLivre}

\rubrique{Problèmes}
\begin{blocLivre}{ex-13-11}\begin{exercice}[3]{Construire sans quadrillage}
Construire l’image du triangle EFG par la translation qui transforme E en K. Décrire la construction et la justifier.
\label{fig-ex-13-11}\figureGeo{trans-11}
\end{exercice}\end{blocLivre}
\begin{blocLivre}{ex-13-12}\begin{exercice}[3]{Translation et parallélogramme}
ABCD est un parallélogramme. On considère la translation qui transforme A en B. Démontrer qu’elle transforme D en C.
\label{fig-ex-13-12}\figureGeo{trans-12}
\end{exercice}\end{blocLivre}
\begin{blocLivre}{ex-13-13}\begin{exercice}[3]{Un pavage par translations}
Un motif carré de 2 cm de côté est répété sans trou ni chevauchement sur 6 colonnes et 4 lignes. Combien de motifs utilise-t-on ? Donner l’aire totale et le glissement qui mène d’un motif au motif situé deux colonnes plus à droite et une ligne plus haut.
\label{fig-ex-13-13}\figureGeo{trans-13}
\end{exercice}\end{blocLivre}
\begin{blocLivre}{ex-13-14}\begin{exercice}[3]{Deux parterres}
Le parterre triangulaire ABC est rectangle en A, avec $AB=6$ m, $AC=8$ m et $BC=10$ m. On le reproduit par la translation qui transforme A en D ; B a pour image E et C pour image F. Quelle longueur de bordure faut-il pour entourer les deux parterres ? Quelle est leur aire totale ? Quelle est la nature de ABED ?
\label{fig-ex-13-14}\figureGeo{trans-14}
\end{exercice}\end{blocLivre}
\begin{blocLivre}{ex-13-15}\begin{exercice}[3]{Un angle transporté}
Dans le triangle ABC, $\widehat{ABC}=70^\circ$. La translation qui transforme B en C transforme A en E. Démontrer que ABCE est un parallélogramme, puis calculer $\widehat{BCE}$.
\label{fig-ex-13-15}\figureGeo{trans-15}
\end{exercice}\end{blocLivre}

\clearpage
